\section{The field selected by the curvature intersection}
\label{V:null-field}

We now construct the base needed for Theorem~\ref{V:exact}. The data
$X,L,\alpha,\nu,G_0,\ell$ remain as in the introduction. A divisor
is called \emph{null} in this section if it is effective and has
$\ell$-value zero. Nonnegativity of $\ell$ implies that every effective
divisor bounded by a null divisor is again null.

\begin{lemma}\label{V:polar-field}
The set
\[
 k=\{0\}\cup\{h\in\C(X)^\times:\ell((h)_\infty)=0\}
\]
is a finitely generated subfield of $\C(X)$, relatively algebraically
closed in $\C(X)$.
\end{lemma}
\begin{proof}
The polar divisors of $h_1+h_2$ and $h_1h_2$ are bounded by the sum
of the polar divisors of $h_1,h_2$. The zero and polar divisors of a
nonzero rational function are linearly equivalent, and therefore have
the same $\ell$-value. These observations prove closure under addition,
multiplication and inversion.

Suppose that $b\in\C(X)$ is algebraic over $k$. In a monic equation
$b^r+a_1b^{r-1}+\cdots+a_r=0$, each coefficient belongs to $k$.
At a prime divisor $P$ where $b$ has a pole, the term $b^r$ cannot
have strictly smaller order than every other term. Thus, for some $i$,
$i\ord_P(b)\ge\ord_P(a_i)$. In particular the pole order of $b$
is bounded by the sum of the pole orders of the coefficients. Its polar
divisor is null, so $b\in k$.

Choose a transcendence basis $t_1,\ldots,t_d$ of $k/\C$.
The algebraic closure of $\C(t_1,\ldots,t_d)$ inside the finitely
generated field $\C(X)$ is a finite extension of that rational function
field. The intermediate field $k$ is algebraic over the same field
and is therefore finite over it. This proves finite generation.
\end{proof}

Choose smooth projective models and morphisms
\[
 X\xleftarrow{\ \pi\ } Z\xrightarrow{\ g\ }Y,
 \qquad \C(Y)=k,
\]
with $\pi$ birational and $g$ surjective. Characteristic zero and relative
algebraic closedness imply that its generic fiber is geometrically
integral. We first verify two properties which persist under every
birational change of these models.

\begin{lemma}\label{V:null-rank}
For every very ample divisor $H$ on $Y$,
\begin{equation}\label{V:ample-null}
 \ell(\pi_*g^*H)=0.
\end{equation}
If $E_Z=K_{Z/X}$ is the effective Jacobian divisor, then
\begin{equation}\label{V:zero-rank}
 \dim_{\C(Y)}H^0(Z_\eta,\OO_Z(E_Z)|_{Z_\eta})=1.
\end{equation}
\end{lemma}
\begin{proof}
Choose a basis $s_0,\ldots,s_N$ of the very ample system. The functions
$s_i/s_0$ belong to $k$, so their polar divisors on $X$ are bounded by
one null divisor $D$. Their poles on $Z$ are then bounded by $\pi^*D$.
Base point freeness says that the componentwise maximum of those polar
divisors on $Z$ is $g^*\divisor(s_0)$. After pushforward, effectivity
and monotonicity give
$0\le\ell(\pi_*g^*H)\le\ell(D)=0$.

The Jacobian section $s_{E_Z}$ gives one nonzero section in
\eqref{V:zero-rank}. Divide any other generic-fiber section by it.
The resulting rational function $a\in\C(Z)=\C(X)$ has horizontal
poles bounded by $E_Z$. There are finitely many remaining polar primes.
Each maps into a proper closed subset of $Y$. A sufficiently high very
ample system has an effective member $B$ vanishing to prescribed orders
on their images; increasing these orders ensures
\[
 (a)_\infty\le E_Z+g^*B.
\]
This also works when a vertical prime maps to codimension greater than
one. Since $E_Z$ is exceptional over $X$, pushforward and
\eqref{V:ample-null} show that the polar divisor of $a$ on $X$ is null.
Thus $a\in k$. Every generic section is a scalar multiple of the
Jacobian section, as claimed.
\end{proof}

The equality in \eqref{V:zero-rank} concerns the full adjoint space in
degree zero. It has not been deduced from an eventual rank statement.
It identifies the integrated Jacobian volume with a line metric on the
base, rather than with a selected subspace of a larger direct image.

\subsection{An equidimensional model without a field extension}
Assume $d=\dim Y>0$. We need a normal projective variety $W$ and a
diagram
\begin{equation}\label{V:model}
 X\xleftarrow{\ p\ }W\xrightarrow{\ f\ }Y
\end{equation}
with $p$ birational, $Y$ smooth projective, and $f$ equidimensional
and toroidal for strict toroidal embeddings. All modifications here
are birational; in particular the field of the base remains $k$.
A toroidal embedding is locally modeled on a toric variety with its
dense torus; the complement is its boundary. Strictness means that the
irreducible boundary components are normal. Toroidal morphisms are
locally toric morphisms, so boundary strata determine rational cones and
the morphism determines their lattice-linear maps.

Weak toroidalization \cite[Theorem~1.1]{ADK2013} first gives a
birational toroidal diagram with smooth projective target. For a
toroidal morphism, equidimensionality is equivalent to the condition
that every source cone maps onto a target cone
\cite[Lemma~4.1 and Remark~4.2]{AK2000}. Refine the target cone
complex so that the image of every source cone is a union of cones.
Projective refinements exist by \cite[Lemma~4.3]{AK2000}; further
projective toroidal resolution makes the target nonsingular. Intersecting
the source cones with inverse images of target cones gives projective
source subdivisions and a compatible toroidal morphism. In a target
cone of the new subdivision, an old cone image meets it in a union of
faces which is convex, hence in a face. Thus the new source cones map
onto target cones. This is the equidimensionalization of
\cite[Proposition~4.4]{AK2000}.

In local toric models, a source stratum maps submersively onto its
target stratum. Cone surjectivity bounds the dimension of the fiber in
each stratum above by $n-d$; the opposite bound for full fibers follows
because the target is smooth. This gives equidimensionality. The normal
source may remain singular, and no covering to reduce multiplicities
is needed. Resolve it by $Z\to W$ when using smooth differential forms;
that resolution need not itself be equidimensional. We continue to write
$\pi:Z\to X$ and $g:Z\to Y$ for the induced maps.

\subsection{The horizontal curvature patch}
The next calculation verifies the remaining analytic hypothesis of the
toroidal-volume theorem. Let $\omega_0$ be a positive form representing
$c_1(G_0)$, and let $\omega_H$ represent a very ample class on $Y$.
When $\nu<n$, equation~\eqref{V:ample-null} implies
\begin{equation}\label{V:vanishing-form}
 (\pi^*\alpha)^\nu\wedge g^*\omega_H
       \wedge\pi^*\omega_0^{n-\nu-1}=0.
\end{equation}
Indeed the left side is a semipositive top-degree form whose integral
is the displayed intersection number. It therefore vanishes pointwise.

At a maximum-rank point away from the Jacobian divisor, diagonalize
$\pi^*\alpha$ in coordinates orthonormal for $\pi^*\omega_0$.
Its first $\nu$ eigenvalues are positive and the rest are zero.
The left side of \eqref{V:vanishing-form} is a positive multiple of
the sum of the diagonal entries of $g^*\omega_H$ in the remaining
$n-\nu$ directions. These entries vanish. A semipositive Hermitian
form vanishing on a vector has zero mixed terms with that vector, so
\[
 \ker\pi^*\alpha\subset\ker dg.
\]
For $\nu=n$ this inclusion is automatic. The maximum-rank locus is a
nonempty ordinary open and intersects the dense smooth base-change
locus, away from the Jacobian divisor. Near a point in their intersection
the rank is constant. On a smaller relatively compact neighborhood the
positive eigenvalues are bounded below, and the kernel inclusion gives
\begin{equation}\label{V:strict-patch}
 \pi^*\alpha\ge c\,g^*\omega_H\qquad(c>0).
\end{equation}
The Jacobian section is nonzero on this patch. If $\nu=0$, the same
calculation would force $dg=0$ on a nonempty open; thus a positive-
dimensional base cannot occur in that case.

We have now proved both analytic inputs from the null-polar field:
full rank one in \eqref{V:zero-rank} and the patch
\eqref{V:strict-patch}. Apply the toroidal-volume theorem
\cite[Theorem \textup{(Toroidal volume metric)}]{CompanionG} to \eqref{V:model}.
For a prime $E$ of $W$ put $a_E=\ord_E K_{W/X}$. Its conclusion uses
\begin{equation}\label{V:correction}
 A_0=\sum_Q\max\left\{0,
       \min_{E\mapsto Q}\left(\frac{1+a_E}{\ord_E f^*Q}-1\right)
       \right\}Q.
\end{equation}
It supplies $A_0\ge0$, $p_*f^*A_0=0$, and a semipositive metric on
$-K_Y+A_0$ with locally integrable squared frame norms, smooth and
strictly positive on a nonempty ordinary open. The cited theorem proves
the all-valuation discrepancy calculation and extension across boundary
divisors and codimension two. Its strictness requirement is exactly
\eqref{V:strict-patch}; no global strict curvature bound is being
inferred directly from that patch.

The correction-ring theorem \cite[Theorem \textup{(Finite generation near a correction divisor)}]{CompanionE}
now applies. For every finite list of rational divisors $P_i$ on $Y$,
there exist $\epsilon\in\Q_{>0}$ and $k_0\in\Z_{>0}$ such that
all $k_0(A_0+\epsilon P_i)$ are Cartier and their multigraded section
ring is finitely generated. This statement retains the actual rational
linear equivalences and the common value of $\epsilon$. We next choose
the two directions dictated by the divisors in Theorem~\ref{V:exact}.
